% 留出验证报告（全量版）
% 数据源：tda_project/outputs/regime_v018_*（2026-11-16 封存开箱产物）
\documentclass[12pt,a4paper]{ctexart}
\usepackage{geometry}
\geometry{left=2.5cm,right=2.5cm,top=2.5cm,bottom=2.4cm}
\usepackage{booktabs}
\usepackage{tabularx}
\usepackage{array}
\usepackage{amsmath,amssymb}
\usepackage{enumitem}
\usepackage{hyperref}
\hypersetup{hidelinks}
\setlength{\parindent}{2em}
\setlength{\parskip}{4pt}
\linespread{1.25}
\setcounter{secnumdepth}{2}
\setcounter{tocdepth}{2}
\renewcommand{\arraystretch}{1.22}
\newcolumntype{Y}{>{\raggedright\arraybackslash}X}

\begin{document}

\begin{titlepage}
\vspace*{2.6cm}
\begin{center}
{\LARGE\bfseries 留出验证报告（全量版）\par}
\vspace{0.7cm}
{\large 按预登记协议检验冻结规则\quad 研究档案 v0.18\par}
\vspace{0.5cm}
{\normalsize 含四层逐事件归因、误报段全表、协议条款与口径换算\par}
\vspace{2.2cm}
{\normalsize \textbf{作者}：潘炯铠\par\vspace{0.35em}}
{\normalsize \textbf{项目}：持续同调量化研究项目（tda\_project）\par\vspace{0.35em}}
{\normalsize \textbf{协议版本}：regime-sealed-validation-v0.18-2026-11-16\par\vspace{0.35em}}
{\normalsize \textbf{登记时刻}：2026-11-16 22:50\quad\textbf{首次计算结果}：2026-11-16 23:09\par\vspace{0.35em}}
{\normalsize \textbf{报告日期}：2026 年 11 月 22 日\par}
\vspace{2.0cm}
\begin{minipage}{0.84\textwidth}
\small
\textbf{摘要}：本报告把训练期（2001-01-10 至 2025-04-07）已经冻结的危机识别与预案规则，原样应用到协议登记的验证段（2025-01-01 至 2026-11-03，442 个交易日），不调参，也不因结果好坏修改规则。检验分别回答四个问题：有没有告警、告警有多早、类型判别是否可靠、误报有多少。六段危机样本均有告警覆盖，其中两段按报告口径提前 20 至 23 个交易日；但两段可评价的类型判断均未判对，其中一段按相反类型执行了预案。严格排除与训练期事件重叠的日期后，误报为 6 日（占 1.6\%），2026 年截至 11 月 3 日为零；观测期仅约 1.5 年，不能据此宣称误报水平已经稳定改善。全部结论限于这段历史，不外推为未来实盘能力。
\end{minipage}
\end{center}
\end{titlepage}

\tableofcontents
\newpage

\section{导言}

\subsection{为什么要有这份报告}

训练期的触发规则、告警管理方式和确认门槛，都是在同一段已知历史上比较后选定的。即使每次比较都诚实执行，最后胜出的方案仍可能沾了“刚好适合这段历史”的便宜。

降低这种选择偏置的一个关键做法，是把选定规则\textbf{冻结}，再把它原样应用到之后的行情。本次检验要回答的不是“这个想法听起来是否合理”，而是“\textbf{这套已经定型的规则，在没有参与选型的行情上表现如何}”。除下文明确披露的日期重叠外，验证数据不参与参数选择。

\subsection{报告结构}

第二章给出协议：登记了什么、什么不能改、留下了哪些修订记录。第三章用直白语言说明系统每天做什么。第四章说明事件归因、时点分档与误报定义。第五章给出完整结果，包括四层逐事件归因与误报段明细。第六章解释误报的两个统计口径。第七章与训练期对照。第八章列出类型短板的事后诊断，明确标注为非预登记、只提供机制线索。第九章列出已知瑕疵。第十章给出复现步骤。

\subsection{阅读时点与证据边界}

本报告保存的是 2026-11-16 对冻结方案一的首次封存检验。后续候选版本曾再次读取这一验证区间，因此那些后续结果只能称为描述性历史重放，不能再算新的独立样本外证据。本次语言润色只改善可读性，不重跑实验，也不改写任何数字、规则或否定结论。

\section{协议}

\subsection{登记信息与纪律}

\begin{center}\small
\begin{tabularx}{\textwidth}{@{}>{\raggedright\arraybackslash}p{3.4cm}Y@{}}
\toprule
\textbf{项} & \textbf{内容} \\
\midrule
协议编号 & regime-sealed-validation-v0.18-2026-11-16 \\
登记时刻 & 2026-11-16 22:50（计算启动前写定并落盘）\\
首次计算结果 & 2026-11-16 23:09（汇总文件生成时刻）\\
验证窗口 & 2025-01-01 至 2026-11-03，442 个交易日 \\
数据截止日 & 2026-11-03；第三段叙事事件称持续至 2026-11-16，后续走势不在数据内 \\
检验目标 & 原样运行训练期冻结规则，核对告警覆盖、时点、判型与误报；不是重新调参 \\
\bottomrule
\end{tabularx}
\end{center}

预登记在计算前明确禁止：调参、更换窗口、重新训练、为追求更好结果修改规则、重判训练期已经冻结的研究结论，以及把本段结果外推为实时能力。

这些限制的目的，是提前锁死“看到结果以后还能改什么”。检验过程中确实出现了几处实现口径需要澄清的问题（见 2.4），每次都有时间和原因记录，并在结果定稿前完成；没有一次是因为结果不好看而修改产品规则。

\subsection{哪些日期真正属于样本外}

协议登记的验证窗口从 2025-01-01 起，但 \textbf{2025-01-01 至 2025-04-07 不是严格的新样本}：事件 16 的标签止于 2025-01-10，事件 17 的事件区间为 2025-03-07 至 03-28、标签延伸至 2025-04-07。这两段都与训练期冻结事件相连。因此，不与训练期事件重叠的验证材料自 \textbf{2025-04-08} 起。

这个区别直接决定误报如何解读：协议口径的 36 个误报日中，有 30 个落在上述重叠区间（见第六章）。所以报告同时给出完整协议窗口和严格样本外窗口，不用一个数字掩盖另一种口径。

\subsection{两类事件分开记账}

验证期有两类事件来源：N1—N3 是按行情叙述划定的“叙事事件”，L1—L3 是按固定标签规则重新计算的“标签事件”。同一波下跌可能同时出现在两类事件中，所以必须分别核对，不能把它们当成六次彼此独立的危机，也不能混算命中率。

\begin{center}\small
\begin{tabularx}{\textwidth}{@{}>{\raggedright\arraybackslash}p{1.1cm}>{\raggedright\arraybackslash}p{4.6cm}>{\raggedright\arraybackslash}p{2.6cm}Y@{}}
\toprule
\textbf{编号} & \textbf{事件} & \textbf{区间} & \textbf{来源与定性} \\
\midrule
N1 & 三月中下旬地缘输入性冲击 & 2026-03-18 至 03-23 & 项目记录的叙事背景（中东冲突、油价冲击），未独立核验；宽基约 $-7\%$，一周修复；对照定性为浅快 \\
N2 & 七月上中旬科技链杠杆踩踏 & 2026-07-09 至 07-31 & 项目记录的叙事背景（存储芯片杠杆破裂、拥挤度极值），未独立核验；宽基约 $-11.5\%$，主线板块深跌；对照定性为浅快（板块深、宽基浅）\\
N3 & 十月下旬起 AI 主线泡沫破裂 & 2026-10-20 至 11-03 & 项目记录的叙事背景（估值回归、红利新高），未独立核验；仍在演化，受数据截止日限制；对照定性为深慢 \\
L1 & 标签段一 & 2026-02-25 至 03-02 & 冻结标签规则（未来 20 个交易日最低累计 $\leq -10\%$）在万得全A 上重算 \\
L2 & 标签段二 & 2026-06-18 至 07-10 & 同上 \\
L3 & 标签段三 & 2026-09-09 至 09-28 & 同上；2026-09-29 之后不足 20 个未来交易日，按规则无法继续标定 \\
\bottomrule
\end{tabularx}
\end{center}

叙事定性（N1 浅快、N2 浅快、N3 深慢）\textbf{只用于事后对照，不作为系统输入}，避免用事后认知反过来修改产品行为。这里的“深慢／浅快”是两套市场形态与预案标签，不是对最终跌幅和持续时间的精确预判。

\subsection{四条实现口径的修订记录}

\begin{center}\small
\begin{tabularx}{\textwidth}{@{}>{\raggedright\arraybackslash}p{1.2cm}>{\raggedright\arraybackslash}p{1.5cm}Y@{}}
\toprule
\textbf{编号} & \textbf{时刻} & \textbf{内容与理由} \\
\midrule
a1 & 23:05 & 首次计算时，同一段告警只能分配给一类事件，导致另一类被误记为漏报。修订后按既定纪律分开归因，同一告警段可在叙事事件与标签事件中各记一次；只有两本账均未归因的才计为误报。首次计算的产物作废，不进入结果。\\
a2 & 23:05 & 首次告警偏移与既有识别窗对齐：同时保留原始值和窗口内统一值，早于认定窗边界的统一记为 $-20$ 个交易日；与训练期中位数比较时使用窗口内统一值。\\
a3 & 23:05 & 误报表增加“是否涉及训练期冻结事件”一列，用于标出与事件 16、17 认定窗的交集；只作注记，不改变判定。\\
a4 & 23:20 & 按登记原文修正误报定义：与两本账任一认定窗有交集的告警段不是误报；有交集但不是该事件首个命中段的，记为“关联命中段”。此前实现误把“未指定为首个命中段”当成误报，相关数字已作废并重算。\\
\bottomrule
\end{tabularx}
\end{center}

\subsection{协议自带的四条已知警示}

预登记在计算前写明四条限制。以下按原意列示，它们直接影响结果解读：

\begin{enumerate}[leftmargin=2em,itemsep=2pt]
\item N3 仍在演化，L3 的后续标签也受数据范围限制；数据止于 2026-11-03。因此，随访不足是\textbf{数据边界}，不能直接记为产品失败。
\item 宽基标签要求未来 20 个交易日最低累计收益不超过 $-10\%$。当宽基只跌 5\% 至 10\%、而主线板块深跌时，标签可能无法覆盖这种结构性危机；产品告警与标签错位时，判定仍以冻结规则为准，行情叙述另行标注。
\item 存储芯片杠杆、地缘冲突、估值传闻等叙事事实来自项目记录，\textbf{未作独立核验}；计算只使用万得全A 基准与冻结产物。
\item 三段标签事件在数据资产审计中先于本预登记算出，已如实披露；其规则与冻结标签逐字相同，没有根据结果调整。
\end{enumerate}

\section{被检验的冻结规则}

\subsection{系统每天收盘后做什么}

系统按顺序做三件事：先检查是否需要提高警觉；如果告警，再检查异常程度和当前市场形态；只有证据充分，才切换风控预案。\textbf{告警不等于立即调仓}。下表中的 S3、B10、REF8、G0 都是研究档案编号，正文中同时给出功能名称，阅读时不必只记代号。

\begin{center}\small
\begin{tabularx}{\textwidth}{@{}>{\raggedright\arraybackslash}p{3.0cm}>{\raggedright\arraybackslash}p{5.2cm}Y@{}}
\toprule
\textbf{环节} & \textbf{冻结规则} & \textbf{本检验中的对照} \\
\midrule
三项条件取二的触发规则（S3） & 三项中至少两项成立才告警：最近 5 日简单收益之和 $\leq -4\%$；20 日波动高于此前 756 日历史的 70 分位；距 20 日高点的回撤 $\geq 5\%$ & 替代触发规则（REF8）\\
五日撤销、十日冷却的告警管理（B10） & 连续 5 日没有原始触发信号便撤销告警；撤销后冷却 10 个交易日，冷却期内不重新告警 & 不做告警状态管理（G0）\\
状态确认与类型判别 & 7 项状态确认指标中至少 4 项达到极端条件，另有 3 项类型指标投票；同一类型连续 2 日成立且第二日确认达标才切换 & 不适用（四个对照层共用同一判定规则）\\
\bottomrule
\end{tabularx}
\end{center}

\subsection{两套风控预案}

切换后按类型选择预案：深慢型采用 50\% 仓位、20\% 止损并买入认沽；浅快型采用 80\% 仓位、12\% 止损。认沽权利金在切换时计入成本，到期赔付按预登记规则进入单位净值路径。这些参数在训练期选定后未再改动。“深慢／浅快”只是市场形态与预案标签，不是对最终跌幅和持续时间的预判；无法判定或数据不足时，不强行选择预案。

\subsection{为什么设置四个对照层}

四层用于两组对照，而不是完整的全因子实验：S3 对 REF8，比较两种触发规则在这段验证行情中的告警与误报；B10 对 G0，比较加入和不加入告警状态管理的差别。最终产品是 S3 加 B10，也是本报告的主结果层。四层共用同一套状态确认与类型判别规则，避免把上游差异和判定差异混在一起。

\section{归因方法}

\subsection{怎样算覆盖，怎样计算首次告警时点}

每个事件都有一个允许告警计入的时间范围，即“事件起点前 20 个交易日至结束后 20 个交易日”，以下简称“认定窗”，与训练期识别窗一致。告警段与认定窗有交集即记为覆盖；多段命中时取最早一段。首次告警日取该段内第一个告警日，首次告警偏移为“首次告警日减事件起点”的交易日数，负值表示提前，正值表示滞后。

时点分四档：首次告警早于起点为“提前”；落在起点与结束之间为“窗内”；落在结束后 20 个交易日内为“结束后”；没有告警为“漏报”。\textbf{覆盖只表示告警段与认定窗相交，不等于提前预警。}

\subsection{为什么偏移要保留两列}

原始偏移可能远低于 $-20$，但认定窗最早只追溯到 $-20$。如果只展示原始值，容易把“提前 46 个交易日”误读成系统可在完整 46 日范围内验证的预警能力。因此，报告同时给出原始值和窗口内统一值：原始偏移低于 $-20$ 时，窗口内统一记为 $-20$，表示“认定窗一开始就已经在告警”。两列并排，用于区分实际日期差与评估窗口边界。

\subsection{误报定义}

一个告警段只有在\textbf{与叙事事件、标签事件的任何认定窗都没有交集}时才计为误报。有交集但不是该事件首个命中段的，记为“关联命中段”，不计入误报。修订 a4 就是为了统一这条定义；此前实现误把“未指定为首个命中段”当成误报，造成误报数虚高，相关数字已作废并重算。

\subsection{类型判别与评价}

叙事类型只用于 N1—N3 的事后对照，标签事件没有独立的叙事类型。表中的系统处理结果分为已评估、未切换和随访不足跳过：未切换表示确认或类型证据没有达到预案切换条件；随访不足表示数据尚未覆盖后续走势，不能评价成败。

“无法判定”在这里是\textbf{产品判定层的结果}，不是对危机识别能力的否定：系统可能已经有告警，但异常程度或类型证据不足，因此不强行选择深慢或浅快。结果表中的“无法判定”字段在未发生确认切换时会保留默认值，也不能被当作逐日类型投票的完整轨迹。

\section{结果}

\subsection{四层总表}

\begin{center}\small\setlength{\tabcolsep}{3.5pt}
\begin{tabularx}{\textwidth}{@{}>{\raggedright\arraybackslash}p{2.9cm}>{\raggedright\arraybackslash}p{1.0cm}>{\raggedright\arraybackslash}p{1.15cm}>{\raggedright\arraybackslash}p{1.0cm}>{\raggedright\arraybackslash}p{1.0cm}>{\raggedright\arraybackslash}p{1.0cm}>{\raggedright\arraybackslash}p{1.0cm}>{\raggedright\arraybackslash}p{1.0cm}>{\raggedright\arraybackslash}p{1.0cm}Y@{}}
\toprule
\textbf{层} & \textbf{告警段} & \textbf{叙事覆盖} & \textbf{标签覆盖} & \textbf{误报段} & \textbf{误报日} & \textbf{关联段} & \textbf{关联日} & \textbf{告警日} & \textbf{训练期同列} \\
\midrule
触发规则（S3，选定） & 13 & 3/3 & 3/3 & 2 & 6 & 6 & 32 & 61 & 9.9 \\
\textbf{告警管理（B10，选定）} & \textbf{8} & \textbf{3/3} & \textbf{3/3} & \textbf{3} & \textbf{36} & \textbf{1} & \textbf{14} & \textbf{113} & \textbf{29.9} \\
替代触发（REF8） & 54 & 3/3 & 3/3 & 37 & 62 & 11 & 54 & 167 & 61.6 \\
不做告警状态管理（G0） & 15 & 3/3 & 3/3 & 4 & 23 & 6 & 32 & 78 & 28.2 \\
\bottomrule
\end{tabularx}
\end{center}

这张表有三点读法。第一，\textbf{四个对照层都覆盖了 6/6}，包括只保留原始触发信号的结果；说明在这段验证行情里，是否告警不是主要瓶颈，类型判别才是。第二，加入告警管理（B10）后的误报日\textbf{多于}原始触发结果（36 对 6），因为五日撤销和十日冷却会改变并可能延长误报段；它的价值是减少反复撤销、重启和实际激活成本，不是减少告警日总数。第三，替代触发规则的误报明显更多（62 对 36 日、37 对 3 段），支持本次验证中选定触发规则的表现；但这只是一段历史，不能外推为未来稳定优势。

\subsection{S3 触发规则：逐事件归因}

% 自动生成：逐事件归因（24 行 = 4 层 x 6 事件）
\subsection*{三项条件取二的触发规则 S3（选定）}
\begin{center}\scriptsize\setlength{\tabcolsep}{3pt}
\begin{tabularx}{\textwidth}{@{}>{\raggedright\arraybackslash}p{0.8cm}>{\raggedright\arraybackslash}p{1.5cm}>{\raggedright\arraybackslash}p{1.5cm}>{\raggedright\arraybackslash}p{1.5cm}>{\raggedright\arraybackslash}p{1.0cm}>{\raggedright\arraybackslash}p{0.85cm}>{\raggedright\arraybackslash}p{0.95cm}>{\raggedright\arraybackslash}p{1.4cm}>{\raggedright\arraybackslash}p{1.15cm}>{\raggedright\arraybackslash}p{1.1cm}Y@{}}
\toprule
\textbf{编号} & \textbf{起点} & \textbf{结束} & \textbf{首次告警日} & \textbf{原始偏移} & \textbf{窗口值} & \textbf{分档} & \textbf{判型} & \textbf{预案} & \textbf{系统状态} & \textbf{说明} \\
\midrule
N1 & 2026-03-18 & 2026-03-23 & 2026-03-19 & 1 & 1 & 窗内 &  &  &  & 窗内第 1 日 \\
N2 & 2026-07-09 & 2026-07-31 & 2026-06-05 & -23 & -20 & 提前 &  &  &  & 提前 20 日 \\
N3 & 2026-10-20 & 2026-11-03 & 2026-09-11 & -20 & -20 & 提前 &  &  &  & 提前 20 日 \\
L1 & 2026-02-25 & 2026-03-02 & 2026-03-19 & 16 & 16 & 结束后 &  &  &  & 窗后 16 日 \\
L2 & 2026-06-18 & 2026-07-10 & 2026-06-01 & -13 & -13 & 提前 &  &  &  & 提前 13 日 \\
L3 & 2026-09-09 & 2026-09-28 & 2026-08-24 & -12 & -12 & 提前 &  &  &  & 提前 12 日 \\
\bottomrule
\end{tabularx}
\end{center}
\subsection*{五日撤销、十日冷却的告警管理 B10（选定）}
\begin{center}\scriptsize\setlength{\tabcolsep}{3pt}
\begin{tabularx}{\textwidth}{@{}>{\raggedright\arraybackslash}p{0.8cm}>{\raggedright\arraybackslash}p{1.5cm}>{\raggedright\arraybackslash}p{1.5cm}>{\raggedright\arraybackslash}p{1.5cm}>{\raggedright\arraybackslash}p{1.0cm}>{\raggedright\arraybackslash}p{0.85cm}>{\raggedright\arraybackslash}p{0.95cm}>{\raggedright\arraybackslash}p{1.4cm}>{\raggedright\arraybackslash}p{1.15cm}>{\raggedright\arraybackslash}p{1.1cm}Y@{}}
\toprule
\textbf{编号} & \textbf{起点} & \textbf{结束} & \textbf{首次告警日} & \textbf{原始偏移} & \textbf{窗口值} & \textbf{分档} & \textbf{判型} & \textbf{预案} & \textbf{系统状态} & \textbf{说明} \\
\midrule
N1 & 2026-03-18 & 2026-03-23 & 2026-03-19 & 1 & 1 & 窗内 & 深慢 & 深慢 & 已评估 & 窗内第 1 日 \\
N2 & 2026-07-09 & 2026-07-31 & 2026-06-01 & -27 & -20 & 提前 & 无法判定 & 未切换 & 已评估 & 提前 20 日 \\
N3 & 2026-10-20 & 2026-11-03 & 2026-09-11 & -20 & -20 & 提前 &  &  & 随访不足 & 数据截止，无法随访 \\
L1 & 2026-02-25 & 2026-03-02 & 2026-03-19 & 16 & 16 & 结束后 & 深慢 & 深慢 & 已评估 & 窗后 16 日 \\
L2 & 2026-06-18 & 2026-07-10 & 2026-06-01 & -13 & -13 & 提前 & 无法判定 & 未切换 & 已评估 & 提前 13 日 \\
L3 & 2026-09-09 & 2026-09-28 & 2026-07-07 & -46 & -20 & 提前 & 无法判定 & 未切换 & 已评估 & 提前 20 日 \\
\bottomrule
\end{tabularx}
\end{center}
\subsection*{替代触发规则 REF8（对照）}
\begin{center}\scriptsize\setlength{\tabcolsep}{3pt}
\begin{tabularx}{\textwidth}{@{}>{\raggedright\arraybackslash}p{0.8cm}>{\raggedright\arraybackslash}p{1.5cm}>{\raggedright\arraybackslash}p{1.5cm}>{\raggedright\arraybackslash}p{1.5cm}>{\raggedright\arraybackslash}p{1.0cm}>{\raggedright\arraybackslash}p{0.85cm}>{\raggedright\arraybackslash}p{0.95cm}>{\raggedright\arraybackslash}p{1.4cm}>{\raggedright\arraybackslash}p{1.15cm}>{\raggedright\arraybackslash}p{1.1cm}Y@{}}
\toprule
\textbf{编号} & \textbf{起点} & \textbf{结束} & \textbf{首次告警日} & \textbf{原始偏移} & \textbf{窗口值} & \textbf{分档} & \textbf{判型} & \textbf{预案} & \textbf{系统状态} & \textbf{说明} \\
\midrule
N1 & 2026-03-18 & 2026-03-23 & 2026-02-13 & -17 & -17 & 提前 &  &  &  & 提前 17 日 \\
N2 & 2026-07-09 & 2026-07-31 & 2026-06-05 & -23 & -20 & 提前 &  &  &  & 提前 20 日 \\
N3 & 2026-10-20 & 2026-11-03 & 2026-09-11 & -20 & -20 & 提前 &  &  &  & 提前 20 日 \\
L1 & 2026-02-25 & 2026-03-02 & 2026-01-26 & -16 & -16 & 提前 &  &  &  & 提前 16 日 \\
L2 & 2026-06-18 & 2026-07-10 & 2026-05-21 & -20 & -20 & 提前 &  &  &  & 提前 20 日 \\
L3 & 2026-09-09 & 2026-09-28 & 2026-07-10 & -43 & -20 & 提前 &  &  &  & 提前 20 日 \\
\bottomrule
\end{tabularx}
\end{center}
\subsection*{不做告警状态管理 G0（对照）}
\begin{center}\scriptsize\setlength{\tabcolsep}{3pt}
\begin{tabularx}{\textwidth}{@{}>{\raggedright\arraybackslash}p{0.8cm}>{\raggedright\arraybackslash}p{1.5cm}>{\raggedright\arraybackslash}p{1.5cm}>{\raggedright\arraybackslash}p{1.5cm}>{\raggedright\arraybackslash}p{1.0cm}>{\raggedright\arraybackslash}p{0.85cm}>{\raggedright\arraybackslash}p{0.95cm}>{\raggedright\arraybackslash}p{1.4cm}>{\raggedright\arraybackslash}p{1.15cm}>{\raggedright\arraybackslash}p{1.1cm}Y@{}}
\toprule
\textbf{编号} & \textbf{起点} & \textbf{结束} & \textbf{首次告警日} & \textbf{原始偏移} & \textbf{窗口值} & \textbf{分档} & \textbf{判型} & \textbf{预案} & \textbf{系统状态} & \textbf{说明} \\
\midrule
N1 & 2026-03-18 & 2026-03-23 & 2026-03-19 & 1 & 1 & 窗内 & 深慢 & 深慢 & 已评估 & 窗内第 1 日 \\
N2 & 2026-07-09 & 2026-07-31 & 2026-06-05 & -23 & -20 & 提前 & 无法判定 & 未切换 & 已评估 & 提前 20 日 \\
N3 & 2026-10-20 & 2026-11-03 & 2026-09-11 & -20 & -20 & 提前 &  &  & 随访不足 & 数据截止，无法随访 \\
L1 & 2026-02-25 & 2026-03-02 & 2026-03-19 & 16 & 16 & 结束后 & 深慢 & 深慢 & 已评估 & 窗后 16 日 \\
L2 & 2026-06-18 & 2026-07-10 & 2026-06-01 & -13 & -13 & 提前 & 无法判定 & 未切换 & 已评估 & 提前 13 日 \\
L3 & 2026-09-09 & 2026-09-28 & 2026-08-24 & -12 & -12 & 提前 & 无法判定 & 未切换 & 已评估 & 提前 12 日 \\
\bottomrule
\end{tabularx}
\end{center}

S3 这一层只生成触发信号，不运行状态确认和类型判别，所以“判型／预案／系统状态”三列为空；它只回答“是否告警、何时告警”。

\subsection{四个事实}

\textbf{事实一：告警覆盖 6/6，但只有一段在事件窗内首次告警。}六段都有告警覆盖，但只有 N1 落在窗内（偏移 $+1$）；N2、N3、L2、L3 记为提前，L1 记为结束后。表面看“提前”较多，但必须拆开来源：N2、N3 已到达 $-20$ 的认定窗边界，L3 的 $-46$ 又受到标签规则前视结构影响。报告采用的可解释提前量是 N2 的 23 个交易日与 N3 的 20 个交易日，不能由 6/6 推出“六段都提前”。

\textbf{事实二：L1 与 N1 是同一波下跌的两个切片。}L1 区间为 2026-02-25 至 03-02，万得全A 在区间内收益为 $+0.21\%$、最大回撤仅 $-0.40\%$；但其后 20 个交易日最低累计收益跌至 $-10.58\%$。也就是说，标签段本身几乎没跌，下跌主要发生在标签段之后。首次告警落在 $+16$ 个交易日的“结束后”档，正是这种时间结构的直接结果，而不是提前预警。

\textbf{事实三：N3 的 $-20$ 是窗口边界，不是系统的能力上限。}N3 从 2026-10-20 开始，而首次告警日在 2026-09-11。按窗口内统一口径记为 $-20$，表示认定窗一开始就已经在告警；窗口之外的提前量不在本次评价范围内。

\textbf{事实四：L3 的提前 46 个交易日不能读作预警。}L3 来自“未来 20 个交易日最低累计收益不超过 $-10\%$”的标签规则，标签起点天然早于后续可见下跌。若把这 46 日直接算作预警能力，会高估系统。

\subsection{类型判别：确定可评价的两段判对零段}

这是本次检验最重要的否定结果：\textbf{告警覆盖不能代替类型判别，更不能代替预案有效性。}

\begin{center}\small
\begin{tabularx}{\textwidth}{@{}>{\raggedright\arraybackslash}p{1.0cm}>{\raggedright\arraybackslash}p{3.6cm}>{\raggedright\arraybackslash}p{3.0cm}Y@{}}
\toprule
\textbf{编号} & \textbf{叙事定性（对照）} & \textbf{实际判定} & \textbf{性质与后果} \\
\midrule
N1 & 浅快（宽基约 $-7\%$，一周修复） & 深慢，并据此切换到 50\% 仓位加认沽 & \textbf{判错且方向相反}。对一周内修复的冲击，这一预案的事后成本高于不切换 \\
N2 & 浅快（板块深、宽基浅） & 无法判定，未切换 & 未形成可执行判断。宽基未破 $-10\%$、主线板块深跌，属于标签盲区 \\
N3 & 深慢 & 无（随访不足跳过） & 数据尚未覆盖后续走势，不评价成败 \\
\bottomrule
\end{tabularx}
\end{center}

N1 与 N2 两段具有事后对照、可以评价，均未判对。N1 的问题比“没有判断”更具体：系统在一周内修复的冲击中，采用了为持续回撤形态设计的 50\% 仓位加认沽预案，事后看成本高于不切换。

训练期同项指标为首次判型一致 10/19、最终预案一致 7/19，本身就不高；验证期结果再次暴露同一短板。两段样本不足以估计稳定的类型准确率，也不能单独证明唯一失败原因，但足以否定“当前判型规则已经可靠”的说法。第八章只提供可能机制的事后诊断。

\subsection{误报段明细}

\begin{center}\small
四层共 46 个误报段。选定的告警管理规则（B10）只有 3 段，见下表加粗部分。
\end{center}

% 自动生成：封存段误报段明细（46 行）
\begin{center}\scriptsize\setlength{\tabcolsep}{3.5pt}\renewcommand{\arraystretch}{1.20}
\begin{tabularx}{\textwidth}{@{}>{\raggedright\arraybackslash}p{1.1cm}>{\raggedright\arraybackslash}p{1.75cm}>{\raggedright\arraybackslash}p{1.75cm}>{\raggedright\arraybackslash}p{1.0cm}>{\raggedright\arraybackslash}p{1.6cm}>{\raggedright\arraybackslash}p{1.4cm}Y@{}}
\toprule
\textbf{层} & \textbf{起点} & \textbf{结束} & \textbf{天数} & \textbf{判型} & \textbf{预案} & \textbf{涉及训练期事件} \\
\midrule
S3 & 2025-01-14 & 2025-01-17 & 4 &  &  & — \\
S3 & 2025-11-21 & 2025-11-24 & 2 &  &  & — \\
\midrule
REF8 & 2025-01-13 & 2025-01-23 & 9 &  &  & — \\
REF8 & 2025-04-29 & 2025-05-07 & 4 &  &  & — \\
REF8 & 2025-05-09 & 2025-05-09 & 1 &  &  & — \\
REF8 & 2025-05-13 & 2025-05-13 & 1 &  &  & — \\
REF8 & 2025-05-15 & 2025-05-16 & 2 &  &  & — \\
REF8 & 2025-06-10 & 2025-06-10 & 1 &  &  & — \\
REF8 & 2025-06-13 & 2025-06-13 & 1 &  &  & — \\
REF8 & 2025-06-26 & 2025-06-26 & 1 &  &  & — \\
REF8 & 2025-06-30 & 2025-06-30 & 1 &  &  & — \\
REF8 & 2025-07-02 & 2025-07-02 & 1 &  &  & — \\
REF8 & 2025-07-04 & 2025-07-07 & 2 &  &  & — \\
REF8 & 2025-07-09 & 2025-07-09 & 1 &  &  & — \\
REF8 & 2025-07-15 & 2025-07-15 & 1 &  &  & — \\
REF8 & 2025-07-23 & 2025-07-23 & 1 &  &  & — \\
REF8 & 2025-07-25 & 2025-07-25 & 1 &  &  & — \\
REF8 & 2025-07-30 & 2025-07-31 & 2 &  &  & — \\
REF8 & 2025-08-07 & 2025-08-08 & 2 &  &  & — \\
REF8 & 2025-08-14 & 2025-08-14 & 1 &  &  & — \\
REF8 & 2025-08-19 & 2025-08-19 & 1 &  &  & — \\
REF8 & 2025-08-21 & 2025-08-21 & 1 &  &  & — \\
REF8 & 2025-08-26 & 2025-08-27 & 2 &  &  & — \\
REF8 & 2025-09-02 & 2025-09-04 & 3 &  &  & — \\
REF8 & 2025-09-12 & 2025-09-12 & 1 &  &  & — \\
REF8 & 2025-09-18 & 2025-09-19 & 2 &  &  & — \\
REF8 & 2025-09-26 & 2025-09-26 & 1 &  &  & — \\
REF8 & 2025-10-10 & 2025-10-13 & 2 &  &  & — \\
REF8 & 2025-10-17 & 2025-10-17 & 1 &  &  & — \\
REF8 & 2025-10-30 & 2025-10-31 & 2 &  &  & — \\
REF8 & 2025-11-21 & 2025-11-24 & 2 &  &  & — \\
REF8 & 2025-11-27 & 2025-11-27 & 1 &  &  & — \\
REF8 & 2025-12-03 & 2025-12-04 & 2 &  &  & — \\
REF8 & 2025-12-23 & 2025-12-23 & 1 &  &  & — \\
REF8 & 2025-12-29 & 2025-12-29 & 1 &  &  & — \\
REF8 & 2026-01-13 & 2026-01-13 & 1 &  &  & — \\
REF8 & 2026-04-23 & 2026-04-24 & 2 &  &  & — \\
REF8 & 2026-05-12 & 2026-05-12 & 1 &  &  & — \\
REF8 & 2026-05-14 & 2026-05-15 & 2 &  &  & — \\
\midrule
B10 & 2025-01-02 & 2025-01-23 & 16 & 深慢 & 深慢 & 事件 16 \\
B10 & 2025-04-07 & 2025-04-24 & 14 & 深慢 & 深慢 & 事件 17 \\
B10 & 2025-11-21 & 2025-11-28 & 6 & 深慢 & 深慢 & — \\
\midrule
G0 & 2025-01-02 & 2025-01-10 & 7 & 深慢 & 深慢 & 事件 16 \\
G0 & 2025-01-14 & 2025-01-17 & 4 & 深慢 & 深慢 & — \\
G0 & 2025-04-07 & 2025-04-18 & 10 & 深慢 & 深慢 & 事件 17 \\
G0 & 2025-11-21 & 2025-11-24 & 2 & 深慢 & 深慢 & — \\
\bottomrule
\end{tabularx}
\end{center}

\subsection{认定窗事实}

\begin{center}\small
\begin{tabularx}{\textwidth}{@{}>{\raggedright\arraybackslash}p{1.0cm}>{\raggedright\arraybackslash}p{3.4cm}>{\raggedright\arraybackslash}p{1.9cm}>{\raggedright\arraybackslash}p{1.7cm}>{\raggedright\arraybackslash}p{1.9cm}Y@{}}
\toprule
\textbf{编号} & \textbf{认定窗} & \textbf{区间收益} & \textbf{窗内最大回撤} & \textbf{未来 20 日最低累计} & \textbf{备注} \\
\midrule
N1 & 2026-03-18 至 03-23 & $-7.27\%$ & $-7.27\%$ & — & 一周内修复，浅快 \\
N2 & 2026-07-09 至 07-31 & $-9.83\%$ & $-11.45\%$ & — & 宽基浅、板块深 \\
N3 & 2026-10-20 至 11-03 & $-6.09\%$ & $-6.69\%$ & — & 数据截止，仍在演化 \\
L1 & 2026-02-25 至 03-02 & $+0.21\%$ & $-0.40\%$ & $-10.58\%$ & 窗内几乎不跌，跌幅在窗后 \\
L2 & 2026-06-18 至 07-10 & $-4.04\%$ & $-6.79\%$ & $-14.65\%$ & 跌幅延续至窗后 \\
L3 & 2026-09-09 至 09-28 & $-4.39\%$ & $-4.92\%$ & $-16.14\%$ & 同上 \\
\bottomrule
\end{tabularx}
\end{center}

这张表解释了一个表面矛盾：L1、L2、L3 的窗内跌幅都不到 $-10\%$，为什么仍被标为危机？因为冻结标签看的是\textbf{随后 20 个交易日的最低累计收益}，不是事件窗内已经发生的跌幅。因此，标签起点可以早于主要下跌，标签事件与按行情叙述划定的事件也不必重合。

\section{误报为什么要报两个口径}

\subsection{三段误报分别落在什么位置}

五日撤销、十日冷却的告警管理规则（B10）在 442 个交易日里产生 3 个误报段、共 36 个误报日。前两段仍与训练期冻结事件的认定窗相交，第三段才位于严格的新样本区间：

\begin{center}\small
\begin{tabularx}{\textwidth}{@{}>{\raggedright\arraybackslash}p{3.4cm}>{\raggedright\arraybackslash}p{1.2cm}>{\raggedright\arraybackslash}p{2.2cm}>{\raggedright\arraybackslash}p{2.0cm}Y@{}}
\toprule
\textbf{误报段} & \textbf{天数} & \textbf{涉及训练期事件} & \textbf{判型／预案} & \textbf{性质} \\
\midrule
2025-01-02 至 01-23 & 16 & 事件 16 & 深慢／深慢 & 与训练期冻结事件认定窗重叠，\textbf{非纯样本外} \\
2025-04-07 至 04-24 & 14 & 事件 17 & 深慢／深慢 & 同上，\textbf{非纯样本外} \\
2025-11-21 至 11-28 & 6 & 无 & 深慢／深慢 & 纯样本外 \\
\bottomrule
\end{tabularx}
\end{center}

\subsection{两个口径}

\begin{center}\small
\begin{tabularx}{\textwidth}{@{}>{\raggedright\arraybackslash}p{2.8cm}>{\raggedright\arraybackslash}p{1.9cm}>{\raggedright\arraybackslash}p{1.4cm}>{\raggedright\arraybackslash}p{1.4cm}>{\raggedright\arraybackslash}p{1.4cm}>{\raggedright\arraybackslash}p{1.2cm}Y@{}}
\toprule
\textbf{口径} & \textbf{起算日} & \textbf{交易日} & \textbf{误报日} & \textbf{日/年} & \textbf{占比} & \textbf{说明} \\
\midrule
协议口径 & 2025-01-01 & 442 & 36 & 19.8 & 8.1\% & 封存登记时的定义，含两段非纯样本外 \\
纯样本外口径 & 2025-04-08 & 381 & 6 & 3.8 & 1.6\% & 剔除冻结事件波及段 \\
训练期同口径 & 2001-01-10 & — & — & 29.9 & 12.0\% & 用于判断是否改善 \\
\bottomrule
\end{tabularx}
\end{center}

两个数字各有用途，但不能混合解释。协议口径的 36 个误报日中，有 30 个（占 83\%）位于非纯样本外区间；若直接拿 19.8 与训练期的 29.9 相比，会得到“下降约三分之一”。若只机械比较严格样本外口径，则对应数字为 29.9 与 3.8，差异约 87\%。这两个百分比使用的分子和分母不同，回答的不是同一个问题。

\subsection{为什么 87\% 不能作为已证实的改善幅度}

3.8 这个数只由 6 个误报日和 381 个交易日（约 1.57 年）得到。样本这么短，一个误报段边界的移动就能明显改变年化值；19.8 那一侧又包含 30 个与训练期事件相连的误报日。因此，任何“误报率已经稳定下降某个百分比”的说法都超出证据。

\begin{quote}
\textbf{可以说的}：严格样本外区间只出现 6 个误报日（2025-11-21 至 11-28），2026 年截至 11 月 3 日为零；这段时期未出现训练期量级的误报。\\
\textbf{不能说的}：误报率相对训练期已经稳定下降了某个百分比。
\end{quote}

\subsection{年化基数的说明}

日/年由误报日数除以窗口年数得到。协议口径的 442 个交易日对应约 1.82 年（因此 36 日折 19.8 日/年、6 日折 3.3 日/年）；纯样本外口径的 381 个交易日对应约 1.57 年（6 日折 3.8 日/年）。占比按交易日计：36/442 为 8.1\%，6/381 为 1.6\%。

\section{与训练期对照}

\begin{center}\small
\begin{tabularx}{\textwidth}{@{}>{\raggedright\arraybackslash}p{4.0cm}>{\raggedright\arraybackslash}p{3.8cm}Y@{}}
\toprule
\textbf{指标} & \textbf{训练期（2001-01-10 至 2025-04-07）} & \textbf{验证期（2025-01-01 至 2026-11-03）} \\
\midrule
危机告警覆盖 & 45/45（宽识别窗） & 6/6 \\
时点分档 & 24 起点前、7 窗内、14 结束后 & 0 起点前、1 窗内、1 结束后、4 提前 \\
实际切换 & 19/45 & 2/6（N1 与 L1，两本账各记一次）\\
误报日/年（选定层） & 29.9 & 19.8（协议口径）／3.8（纯样本外）\\
首次判型与评价标签一致 & 10/19 & 0/2 \\
最终预案与评价标签一致 & 7/19 & 0/2 \\
\bottomrule
\end{tabularx}
\end{center}

读这张表要注意三件事。

第一，两边样本量相差很大（45 段对 6 段）。验证期的任何比例都很不稳定，不能当成未来表现的精确参数。

第二，训练期 45/45 与验证期 6/6 都采用“起点前 20 日至结束后 20 日”的宽识别窗，因此“是否覆盖”可以对照；但“提前多少”不能直接比较。训练期有 24 段在起点前触发，而验证期的 N2、N3 已到达认定窗边界，窗口之外的提前量没有进入同一把尺子。

第三，类型判别是本次最明确的负面结果：训练期与评价标签的一致数就不高（10/19、7/19），验证期两段均未判对。两段样本不能分离抽样波动与规则缺陷，但无论具体成因如何，都不能宣称当前类型判别已经可靠。第八章只用于寻找可能机制。

\section{类型短板可能卡在哪里（事后诊断，非预登记）}

本节不属于 v0.18 的预登记检验，而是在看到结果后追加的诊断。\textbf{它只提供机制线索，不作为独立验证结论，也不能直接用来修改规则。}保留这一部分，是为了说明后续应该检查什么，而不是把失败原因写成已经证实的因果链。

\textbf{观察一：2026 年三段样本中，状态确认的整体强度始终不够。}方案一要求 7 项状态确认指标中至少 4 项达到极端条件。在 2026 年三段危机窗内，逐日统计的达标项数\textbf{没有一天达到 4 项}，多数日期只有 1 至 3 项。因此，N2、N3 的“无法判定”主要不是几项指标意见不一，而是证据强度根本没有达到切换门槛。

\textbf{观察二：受限的是“相对自身历史是否极端”，不是类型指标本身。}部分确认指标要与此前的 756 个交易日比较。2026 年这些指标的历史分位大多位于 30 至 70 之间，按规则看起来并不极端。2024 年初微盘股急跌留下的极值又抬高了滚动比较基准，使此后相近的绝对波动水平更难达到极端分位。

\textbf{观察三：行业相关水平处于历史低位。}以 19 个申万一级行业计算，2026 年 20 日平均两两相关的中位数为 0.467，全样本中位数为 0.727；0.467 约处于全样本的 8\% 分位。依赖“相关性上升”的确认指标因此很难达标。在主线板块深跌、其他板块没有同步下跌的结构行情中，行业相关性可以下降，这与“所有危机都必然一起跌”的直觉不同。

三条观察合在一起，提示\textbf{确认门槛与历史比较基准可能不适合这类结构性行情}；但它们不是唯一原因，也没有经过新的前瞻样本验证。后续可以把“按市场状态调整确认条件”或“为低相关结构行情增加独立确认通道”作为新的研究假设，另行预登记检验，不能把本节的诊断直接当成产品升级依据。

\section{已知瑕疵}

\subsection{旧基准日历漏了 19 个交易日}

本次封存检验使用的旧版基准序列有 6,483 个交易日。事后核查发现，它漏掉了 19 个交易日，完整序列应有 6,502 个交易日，并已与独立交易日历核对。

在本报告的 442 日窗口内，漏掉的只有 \textbf{2026-10-15} 一天，当日万得全A 下跌 2.28\%。

该日不在任何事件认定窗内（N3 从 2026-10-20 开始，L1—L3 均在 2026 年 9 月及以前），所以 6/6 覆盖、时点分档与类型结论不受直接影响。但补入该日后，5 日与 20 日滚动窗口会略有移动，理论上可能改变个别误报段的边界。

\textbf{本报告没有按修正后的日历重新计算}，这是必须保留的口径缺口。若要闭合，应另立版本、重新预登记并重跑，不能把现有数字冒充为修正日历后的结果。

\subsection{其余三条}

\begin{enumerate}[leftmargin=2em,itemsep=2pt]
\item N3 与 L3 受数据截止日 2026-11-03 限制，随访不足属于数据边界。
\item 叙事背景来自项目记录，未作独立核验；计算只使用万得全A 与冻结产物。
\item 标签段三段先于本预登记算出，已披露；规则与冻结标签逐字相同。
\end{enumerate}

\section{复现}

\subsection{输入}

预登记声明的只读输入为训练期冻结产物：触发层误报明细、治理层误报明细、基准日线与危机事件表。本报告的四份数据产物均已随附于 \texttt{data} 目录。

\subsection{步骤}

\begin{enumerate}[leftmargin=2em,itemsep=2pt]
\item 确认输入文件的校验值与本报告附录 B 一致。
\item 运行 \texttt{build\_v018\_sealed\_validation.py}，生成汇总结果与三张明细表。
\item 核对 \texttt{regime\_v018\_summary.json} 中的两个校验值：归因表 61799f5cbf82a828、误报表 fc0f817e5b39f053。
\item 若校验值不一致，说明输入或代码版本与登记时不同，结果不可与本报告比较。
\end{enumerate}

\appendix

\section{术语}

\begin{center}\small
\begin{tabularx}{\textwidth}{@{}>{\raggedright\arraybackslash}p{3.6cm}Y@{}}
\toprule
\textbf{术语} & \textbf{含义} \\
\midrule
认定窗 & 允许告警计入某个事件的时间范围：起点前 20 个交易日至结束后 20 个交易日 \\
归属窗 & 用于损益核算起点的时间范围：起点前 20 个交易日至事件结束 \\
首次告警 & 被计入某事件的告警段中的第一个告警日 \\
窗口内统一偏移 & 原始偏移低于 $-20$ 时统一记为 $-20$，表示认定窗一开始就已经在告警 \\
关联命中段 & 与认定窗有交集、但不是该事件首个命中段的告警段，不计入误报 \\
误报段 & 与叙事事件、标签事件的任何认定窗都没有交集的告警段 \\
两本账 & 叙事事件（N1 至 N3）与标签事件（L1 至 L3）分别归因，不混算命中率 \\
无法判定 & 产品的确认或类型证据不足，未强行选择深慢或浅快；不等于漏报 \\
\bottomrule
\end{tabularx}
\end{center}

\section{产物与校验值}

\begin{center}\footnotesize
\begin{tabularx}{\textwidth}{@{}>{\raggedright\arraybackslash}p{6.6cm}Y@{}}
\toprule
\textbf{文件} & \textbf{内容} \\
\midrule
regime\_\hspace{0pt}sealed\_\hspace{0pt}validation\_\hspace{0pt}v0.18.toml & 预登记协议全文，含 a1 至 a4 四条修订 \\
regime\_\hspace{0pt}v018\_\hspace{0pt}summary.json & 四层汇总；归因校验 61799f5cbf82a828，误报校验 fc0f817e5b39f053 \\
regime\_\hspace{0pt}v018\_\hspace{0pt}episode\_\hspace{0pt}attribution.csv & 逐事件归因 24 行（4 层 $\times$ 6 事件）\\
regime\_\hspace{0pt}v018\_\hspace{0pt}false\_\hspace{0pt}alarm\_\hspace{0pt}sealed.csv & 误报段明细 46 行，含冻结事件波及列 \\
regime\_\hspace{0pt}v018\_\hspace{0pt}benchmark\_\hspace{0pt}episode\_\hspace{0pt}facts.csv & 六段认定窗的区间收益、最大回撤、未来 20 日最低 \\
build\_\hspace{0pt}v018\_\hspace{0pt}sealed\_\hspace{0pt}validation.py & 生成脚本 \\
\bottomrule
\end{tabularx}
\end{center}

\vspace{1.5em}
\begin{center}\small
\rule{0.62\textwidth}{0.4pt}\\[0.5em]
本报告记录冻结规则在历史验证段上的表现，不构成对未来行情的预测，\par
亦不构成实盘可用性或收益承诺。
\end{center}

\end{document}
